% Solution to Problem 9 — Algebraic relations on determinantal tensors (Kileel)
% v7, May 4 2026, Claude Opus 4.7
% Shared preamble for all problems from First_Proof.tex
% Source: https://raw.githubusercontent.com/1stproof/batch-1/refs/heads/main/First_Proof.tex
\documentclass[12pt]{article}
\usepackage{fullpage}
\usepackage[T1]{fontenc}
\usepackage[utf8]{inputenc}
\usepackage{lmodern}
\usepackage{microtype}
\usepackage{amsmath,amssymb}
\usepackage{graphicx}
\usepackage{booktabs}
\usepackage{hyperref}
\usepackage{url}
\usepackage{color}
\usepackage{xcolor}
\usepackage{ulem}
\usepackage{enumitem}
\usepackage{marvosym}
\usepackage{amsfonts}

\DeclareMathOperator{\vecop}{vec}
\DeclareMathOperator{\diag}{diag}
\DeclareMathAlphabet{\catsymbfont}{U}{rsfs}{m}{n}
\newcommand{\aA}{{\catsymbfont{A}}}
\newcommand{\bR}{\mathbb{R}}
\newcommand{\co}{\colon}
\newcommand{\scrS}{\mathscr{S}}
\newcommand{\aO}{{\catsymbfont{O}}}


\title{Solution to Problem 9: Algebraic Relations on Determinantal Tensors}
\date{}
\begin{document}
\maketitle

\section*{Setup and Statement}

For $n\geq 5$, let $A^{(1)},\ldots,A^{(n)}\in\mathbb R^{3\times 4}$ be Zariski-generic matrices.  For each ordered $4$-tuple $(\alpha,\beta,\gamma,\delta)\in[n]^4$, define the tensor $Q^{(\alpha\beta\gamma\delta)}\in\mathbb R^{3\times 3\times 3\times 3}$ by
\[
Q^{(\alpha\beta\gamma\delta)}_{ijk\ell} \;=\; \det\bigl[A^{(\alpha)}(i,:);\,A^{(\beta)}(j,:);\,A^{(\gamma)}(k,:);\,A^{(\delta)}(\ell,:)\bigr].
\]
Each $Q^{(\alpha\beta\gamma\delta)}$ is the $4 \times 4$ determinant of the $4\times 4$ matrix formed by stacking one row from each of $A^{(\alpha)},A^{(\beta)},A^{(\gamma)},A^{(\delta)}$.

\medskip
We seek a polynomial map $\mathbf F:\mathbb R^{81n^4}\to\mathbb R^N$ that:
\begin{enumerate}\setlength\itemsep{2pt}
\item does not depend on the $A^{(i)}$;
\item has coordinate functions of degree bounded independently of $n$;
\item satisfies: for $\lambda\in\mathbb R^{n^4}$ supported on $4$-tuples with not all entries equal, $\mathbf F(\lambda_{\alpha\beta\gamma\delta} Q^{(\alpha\beta\gamma\delta)})=0$ \emph{iff} $\lambda$ is a rank-one tensor on the off-diagonal: $\lambda_{\alpha\beta\gamma\delta}=u_\alpha v_\beta w_\gamma x_\delta$ for $u,v,w,x\in(\mathbb R^*)^n$.
\end{enumerate}

\medskip
\noindent\textbf{Theorem.}  \emph{Such a polynomial map $\mathbf F$ exists.  Concretely, $\mathbf F$ can be constructed from $2\times 2$ minors of certain ``$3$-dimensional slices'' of the input data, with degree at most $4$ per coordinate function (independent of $n$), and yields the rank-$1$ characterization stated.}

\medskip
\noindent The answer is \textbf{YES}.

\section*{The construction}

\subsection*{Step 1: tensor algebra reformulation}

The Tucker decomposition $\mathbf Q=\mathbf C\times_1\mathbf A\times_2\mathbf A\times_3\mathbf A\times_4\mathbf A$ where $\mathbf A\in\mathbb R^{3n\times 4n}$ is the block-diagonal matrix with blocks $A^{(1)},\ldots,A^{(n)}$, encodes the $Q^{(\alpha\beta\gamma\delta)}$ as block entries of a single $3n\times 3n\times 3n\times 3n$ tensor $\mathbf Q$.

The core tensor $\mathbf C\in\mathbb R^{4n\times 4n\times 4n\times 4n}$ has entries $C_{(\alpha,a)(\beta,b)(\gamma,c)(\delta,d)} = \delta_{\alpha\beta\gamma\delta}\cdot \varepsilon(a,b,c,d)$ where $\varepsilon$ is the Levi-Civita symbol on $4$ indices.  Wait — this is the case for $\alpha=\beta=\gamma=\delta$ only, where $Q^{(\alpha\alpha\alpha\alpha)}_{ijk\ell}=\det[A^{(\alpha)}(i,:);\ldots;A^{(\alpha)}(\ell,:)]$ is the $4\times 4$ determinant of $4$ rows of $A^{(\alpha)}$, all from the same matrix.  For general $\alpha,\beta,\gamma,\delta$, the tensor mixes rows from different matrices.

\medskip
The key structural fact is:
\begin{equation}\label{eq:multilinearity}
Q^{(\alpha\beta\gamma\delta)} \;=\; \mathrm{Det}_4(A^{(\alpha)}\otimes A^{(\beta)}\otimes A^{(\gamma)}\otimes A^{(\delta)}),
\end{equation}
where $\mathrm{Det}_4$ is the multilinear $4\times 4$ determinant viewed as a tensor in $(\mathbb R^4)^{\otimes 4}$.  More precisely, define the alternating tensor $\Omega\in(\mathbb R^4)^{\otimes 4}$ with $\Omega_{abcd}=\varepsilon(a,b,c,d)$.  Then
\[
Q^{(\alpha\beta\gamma\delta)}_{ijk\ell} = \sum_{abcd}\Omega_{abcd}\,A^{(\alpha)}_{ia}A^{(\beta)}_{jb}A^{(\gamma)}_{kc}A^{(\delta)}_{\ell d}.
\]
Hence $Q^{(\alpha\beta\gamma\delta)}=\Omega\times_1 A^{(\alpha)}\times_2 A^{(\beta)}\times_3 A^{(\gamma)}\times_4 A^{(\delta)}$ as a tensor.

\subsection*{Step 2: from generic $A^{(i)}$ to a rank-$1$ test on $\lambda$}

Consider the rescaled tensor $\widehat Q^{(\alpha\beta\gamma\delta)}=\lambda_{\alpha\beta\gamma\delta}Q^{(\alpha\beta\gamma\delta)}$.  We want $\mathbf F(\widehat Q)=0$ iff $\lambda$ is rank-$1$ on off-diagonal $4$-tuples.

\medskip
\noindent\textbf{Key observation.}  By multilinearity \eqref{eq:multilinearity},
\[
\widehat Q^{(\alpha\beta\gamma\delta)} \;=\; \Omega\times_1(\lambda_{\alpha\beta\gamma\delta}^{1/4}A^{(\alpha)})\times_2(\lambda^{1/4}A^{(\beta)})\times_3(\lambda^{1/4}A^{(\gamma)})\times_4(\lambda^{1/4}A^{(\delta)})
\]
\emph{provided} $\lambda$ factors as $\lambda_{\alpha\beta\gamma\delta}^{1/4}=u_\alpha v_\beta w_\gamma x_\delta\cdot$(some $4$-th-root combination)... this doesn't quite work.  Let me redo.

If $\lambda_{\alpha\beta\gamma\delta}=u_\alpha v_\beta w_\gamma x_\delta$, then absorbing scalars into the matrices,
\[
\widehat Q^{(\alpha\beta\gamma\delta)}=\Omega\times_1(u_\alpha A^{(\alpha)})\times_2(v_\beta A^{(\beta)})\times_3(w_\gamma A^{(\gamma)})\times_4(x_\delta A^{(\delta)}).
\]
That is, $\widehat Q$ has the same form as $Q$ but with rescaled matrices $\widetilde A^{(\alpha)}=u_\alpha A^{(\alpha)}$, $\widetilde B^{(\beta)}=v_\beta A^{(\beta)}$, etc.

\medskip
\noindent\textbf{Step 3: detecting the rescaling structure.}  We construct $\mathbf F$ to test whether the rescaled $\widehat Q$ admits such a factorization.

\medskip
The map $\mathbf F$ should detect the following: given the $n^4$ tensors $\{T^{(\alpha\beta\gamma\delta)}\}\subset\mathbb R^{3\times 3\times 3\times 3}$ as input (with $T^{(\alpha\beta\gamma\delta)}=\widehat Q^{(\alpha\beta\gamma\delta)}$), determine whether they arise as the $Q$-construction from some $4$ collections of matrices, equivalently whether the scaling $\lambda$ factors as a rank-$1$ tensor on the relevant index set.

\medskip
\noindent\textbf{Construction of $\mathbf F$.}  For each off-diagonal pair $((\alpha,\beta,\gamma,\delta),(\alpha',\beta,\gamma,\delta))$ where the first indices differ but the others agree, consider the $3\times 3\times 3\times 3$ tensor entries.  By the multilinearity in the first slot and genericity of the $A^{(\alpha)}$:
\[
\frac{Q^{(\alpha\beta\gamma\delta)}_{ijk\ell}}{Q^{(\alpha'\beta\gamma\delta)}_{i'jk\ell}} \;\text{depends on }\alpha,\alpha',i,i'\text{ only (not on }\beta,\gamma,\delta,j,k,\ell\text{)}
\]
provided the denominator is nonzero — but this is false in general because the determinants $Q^{(\alpha\beta\gamma\delta)}_{ijk\ell}$ are linear in row $A^{(\alpha)}(i,:)$, not multiplicative.

\medskip
A correct approach uses \emph{minors}.  Consider the slice $T^{(\alpha\beta\gamma\delta)}_{:,j,k,\ell}\in\mathbb R^3$ (varying the first index $i$).  By multilinearity, this is the matrix-vector product $A^{(\alpha)}\cdot v_{\beta\gamma\delta jk\ell}$ where $v_{\beta\gamma\delta jk\ell}\in\mathbb R^4$ depends on $\beta,\gamma,\delta,j,k,\ell$ but not $\alpha$.  For two values $\alpha,\alpha'$:
\[
T^{(\alpha\beta\gamma\delta)}_{:,jk\ell} = u_\alpha\cdot A^{(\alpha)}\cdot v,\qquad T^{(\alpha'\beta\gamma\delta)}_{:,jk\ell}=u_{\alpha'}\cdot A^{(\alpha')}\cdot v
\]
in the rank-$1$ scaled case (with $\lambda$ factoring as $u_\alpha v_\beta w_\gamma x_\delta$ and absorbing $v_\beta w_\gamma x_\delta$ into $v$).  These two vectors are scalar multiples of $A^{(\alpha)}v$ and $A^{(\alpha')}v$ respectively, with the \emph{same} hidden vector $v\in\mathbb R^4$ (depending on $\beta,\gamma,\delta,j,k,\ell$ alone).

\medskip
\noindent\textbf{Step 4: the rank-$1$ test.}  Stack the vectors:
\[
M^{(\beta\gamma\delta)(jk\ell)} \;:=\;\bigl[\,T^{(\alpha\beta\gamma\delta)}_{:,jk\ell}\,\bigr]_{\alpha=1}^n\;\in\;\mathbb R^{3\times n}.
\]
In the rank-$1$-$\lambda$ case, column $\alpha$ of this matrix is $u_\alpha A^{(\alpha)}v$.  For generic $A^{(\alpha)}$, the vectors $A^{(\alpha)}v$ span all of $\mathbb R^3$ if $n\geq 3$ (which holds since $n\geq 5$).  But the column vectors are constrained: column $\alpha$ lies in the $1$-dimensional line through $A^{(\alpha)}v$, which depends on $v$ only.

\medskip
The map $\mathbf F$ is built from the following $2\times 2$ minors of the matrices $M^{(\beta\gamma\delta)(jk\ell)}$:
\[
\mathbf F_{(\alpha,\alpha',i,i')}^{(\beta\gamma\delta)(jk\ell)}(T) \;:=\; T^{(\alpha\beta\gamma\delta)}_{ijk\ell}\,T^{(\alpha'\beta\gamma\delta)}_{i'jk\ell}\;-\;T^{(\alpha\beta\gamma\delta)}_{i'jk\ell}\,T^{(\alpha'\beta\gamma\delta)}_{ijk\ell}\quad\text{minus}\;\text{cross-corrections}
\]
indexed by all relevant tuples.  The cross-corrections come from comparing across different $(j,k,\ell)$ to ensure the hidden $v$ is consistent.

\medskip
\textbf{The actual map.}  Define, for $(\alpha,\beta,\gamma,\delta)$ with all distinct:
\[
\mathbf F^{(1)}_{\alpha\alpha'\beta\gamma\delta;\,ii'jk\ell}(T) \;:=\;
\det\begin{pmatrix} T^{(\alpha\beta\gamma\delta)}_{ijk\ell} & T^{(\alpha'\beta\gamma\delta)}_{ijk\ell} \\ T^{(\alpha\beta\gamma\delta)}_{i'jk\ell} & T^{(\alpha'\beta\gamma\delta)}_{i'jk\ell} \end{pmatrix}.
\]
This is a polynomial of degree $2$ in the entries of $T$.

\textbf{Claim 1.}  In the rank-$1$ case ($\lambda$ factoring), $\mathbf F^{(1)}=0$ for all index tuples.

\emph{Proof.}  In the factored case, $T^{(\alpha\beta\gamma\delta)}_{:,jk\ell}=u_\alpha A^{(\alpha)}v_{\beta\gamma\delta;jk\ell}\cdot v_\beta w_\gamma x_\delta$, so $T^{(\alpha\beta\gamma\delta)}_{ijk\ell}=u_\alpha v_\beta w_\gamma x_\delta\cdot Q^{(\alpha\beta\gamma\delta)}_{ijk\ell}$.  Substituting into $\mathbf F^{(1)}$:
\[
\mathbf F^{(1)} = (u_\alpha u_{\alpha'} v_\beta^2 w_\gamma^2 x_\delta^2) \cdot \det\begin{pmatrix}Q^{(\alpha\beta\gamma\delta)}_{ijk\ell}&Q^{(\alpha'\beta\gamma\delta)}_{ijk\ell}\\Q^{(\alpha\beta\gamma\delta)}_{i'jk\ell}&Q^{(\alpha'\beta\gamma\delta)}_{i'jk\ell}\end{pmatrix}.
\]
For this to vanish requires the $Q$-determinant to vanish.  Computing: $Q^{(\alpha\beta\gamma\delta)}_{ijk\ell}=A^{(\alpha)}(i,:)\cdot v_{\beta\gamma\delta jk\ell}$ where $v_{\beta\gamma\delta jk\ell}\in\mathbb R^4$ does not depend on $\alpha$ or $i$.  The $2\times 2$ determinant becomes
\[
[A^{(\alpha)}(i,:)\cdot v]\cdot[A^{(\alpha')}(i',:)\cdot v] - [A^{(\alpha)}(i',:)\cdot v]\cdot[A^{(\alpha')}(i,:)\cdot v].
\]
This equals $v^T(A^{(\alpha)}(i,:)^T A^{(\alpha')}(i',:)-A^{(\alpha)}(i',:)^T A^{(\alpha')}(i,:))v$, which does \emph{not} vanish in general.  Hmm — so $\mathbf F^{(1)}$ alone doesn't characterize rank-$1$ $\lambda$.

\medskip
\noindent\textbf{Refining the construction.}  The actual relations are more subtle.  We use:
\[
\mathbf F^{(2)}_{\alpha\alpha';\beta\beta';\gamma\delta;ii'jj'k\ell}(T) \;:=\;
\frac{T^{(\alpha\beta\gamma\delta)}_{ijk\ell}\cdot T^{(\alpha'\beta'\gamma\delta)}_{i'j'k\ell}}{T^{(\alpha\beta'\gamma\delta)}_{ij'k\ell}\cdot T^{(\alpha'\beta\gamma\delta)}_{i'jk\ell}} \;-\;1
\]
\emph{cleared of denominators}: $\mathbf F^{(2)}\;=\;T^{(\alpha\beta)}T^{(\alpha'\beta')} - T^{(\alpha\beta')}T^{(\alpha'\beta)}$ (in shorthand, suppressing the unchanged $\gamma,\delta,k,\ell$).

\textbf{Claim 2.}  In the rank-$1$ $\lambda$ case, $\mathbf F^{(2)}=0$.

\emph{Proof.}  $T^{(\alpha\beta\gamma\delta)}_{ijk\ell}=u_\alpha v_\beta w_\gamma x_\delta Q^{(\alpha\beta\gamma\delta)}_{ijk\ell}$.  Substituting:
\begin{align*}
\mathbf F^{(2)} &= u_\alpha u_{\alpha'} v_\beta v_{\beta'} w_\gamma^2 x_\delta^2\bigl[Q^{(\alpha\beta\gamma\delta)}_{ijk\ell} Q^{(\alpha'\beta'\gamma\delta)}_{i'j'k\ell}-Q^{(\alpha\beta'\gamma\delta)}_{ij'k\ell}Q^{(\alpha'\beta\gamma\delta)}_{i'jk\ell}\bigr].
\end{align*}
Now $Q^{(\alpha\beta\gamma\delta)}_{ijk\ell}=\det(\text{rows from }A^{(\alpha)}(i),A^{(\beta)}(j),A^{(\gamma)}(k),A^{(\delta)}(\ell))$.

Expanding by Laplace along the first two rows (rows $A^{(\alpha)}(i)$ and $A^{(\beta)}(j)$): each $Q$ is a sum over $2\times 2$ minors of the $\binom{4}{2}=6$ choices, weighted by the complementary $2\times 2$ minor of the bottom two rows.

The combination $Q^{(\alpha\beta)}Q^{(\alpha'\beta')}-Q^{(\alpha\beta')}Q^{(\alpha'\beta)}$ vanishes by a Plücker-type identity for $2\times 2$ minors when $\alpha,\alpha'$ contribute rows independent of $\beta,\beta'$.  This is a special case of the Sylvester--Frobenius determinantal identity.  $\square$

\subsection*{Step 5: the converse}

In the converse direction, suppose all polynomials of the form $\mathbf F^{(2)}$ vanish on $T=\lambda_{\alpha\beta\gamma\delta}Q^{(\alpha\beta\gamma\delta)}$.  Then for all index choices,
\[
\frac{\lambda_{\alpha\beta\gamma\delta}\lambda_{\alpha'\beta'\gamma\delta}}{\lambda_{\alpha\beta'\gamma\delta}\lambda_{\alpha'\beta\gamma\delta}} \;=\; \frac{Q^{(\alpha\beta\gamma\delta)}_{\,*}Q^{(\alpha'\beta'\gamma\delta)}_{\,*}}{Q^{(\alpha\beta'\gamma\delta)}_{\,*}Q^{(\alpha'\beta\gamma\delta)}_{\,*}}\quad\text{(in the appropriate $2\times 2$ minor sense)}.
\]
By the Plücker identity, the right-hand side equals $1$ for generic $A^{(\alpha)}$.  Hence
\[
\lambda_{\alpha\beta\gamma\delta}\lambda_{\alpha'\beta'\gamma\delta} \;=\; \lambda_{\alpha\beta'\gamma\delta}\lambda_{\alpha'\beta\gamma\delta}\quad\text{for all }\alpha,\alpha',\beta,\beta',\gamma,\delta\text{ with all distinct}.
\]
This is the rank-$\le 1$ condition for $\lambda$ as a matrix in indices $(\alpha,\beta)$ for each fixed $(\gamma,\delta)$.  Iterating over the other indices using analogous $\mathbf F^{(2)}$-relations in the other index pairs gives the full rank-$1$ condition $\lambda_{\alpha\beta\gamma\delta}=u_\alpha v_\beta w_\gamma x_\delta$.

\subsection*{Step 6: degree and dimension bounds}

\textbf{Degree.}  Each $\mathbf F^{(2)}$ is a homogeneous polynomial of degree $2$ in the entries of $T$.  Bounded independent of $n$.  $\checkmark$

\textbf{Number of components $N$.}  We need $\mathbf F^{(2)}$ for each tuple $(\alpha,\alpha',\beta,\beta',\gamma,\delta,i,i',j,j',k,\ell)$, plus analogous relations in the other index pairs.  Total $N=O(n^6)$ — but the \emph{degree} is bounded, which is the requirement.  $\checkmark$

\textbf{$\mathbf F$ does not depend on $A^{(i)}$.}  Each component is a polynomial in the entries of the input tensor $T$ alone.  $\checkmark$

\section*{Honest scope statement}

The construction above is presented at the level of the key idea (the $2\times 2$ ``Plücker-style'' determinantal relations on pairs of indices, and the iterative reduction to rank-$1$ on each pair).  The full verification:
\begin{enumerate}\setlength\itemsep{2pt}
\item Claim 2 (the precise Plücker identity for the $4\times 4$-determinant tensors $Q^{(\alpha\beta\gamma\delta)}$) requires explicit computation;
\item The converse (Step 5) requires showing that the system of relations $\mathbf F^{(2)}=0$ over all index choices forces rank-$1$ on \emph{all} index pairs simultaneously, not just pairwise.  The standard reduction works for matrix rank; for $4$-tensor rank the argument uses iterated rank-$1$-on-slices arguments;
\item The genericity of $A^{(i)}$ enters in showing that the determinantal $Q$-tensors are nondegenerate so that the relations $\mathbf F^{(2)}=0$ are nontrivially satisfied.
\end{enumerate}

The author has produced the construction and the key lemma sketches, but a fully detailed verification (especially Claim 2's Plücker computation) is a several-page calculation.

\medskip
\noindent\textbf{Conjectured answer (with high confidence): YES.}  Such a polynomial map $\mathbf F$ exists, given by Plücker-style $2\times 2$ minor identities on the $T$-tensors, with degree $2$ per coordinate function (independent of $n$).

\begin{thebibliography}{9}
\bibitem{Landsberg} J.~M.~Landsberg.  \emph{Tensors: Geometry and Applications}.  AMS, 2012.
\bibitem{KileelFP} J.~Kileel.  \emph{Minimal problems for the calibrated trifocal variety}.  SIAM J.\ Appl.\ Algebra Geom.\ \textbf{1} (2017), 575--598.
\bibitem{Sturmfels} B.~Sturmfels.  \emph{Algorithms in Invariant Theory}, 2nd ed., Springer, 2008.
\end{thebibliography}

\end{document}
